\section{Discussion}

\subsection{Persistent latent structure}

The proposed construction exploits persistent latent structure. A confidence-band robustification asks which null distributions remain compatible with the contaminated reference data and thereby forgets which individual reference observations might be contaminated. Static coupling retains that uncertainty as one fixed latent candidate set
\[
  C\subseteq\{1,\ldots,K\}
\]
that is reused at every future time.

The distinction is therefore not simply \emph{rank method versus CDF method}. It is a distinction between \emph{persistent latent identity uncertainty} and \emph{generic distributional uncertainty}.
The lower envelope can still be conservative, but it preserves the constraint that one contamination pattern must explain the entire online rank history.

\subsection{Relation to the oracle}

At the primary frozen setting
\[
  K=32,\qquad m=2,\qquad \gamma=4,\qquad T=200,
\]
the oracle crossing probability is \(0.9810\). Static coupling reaches \(0.9144\) in the original outside-support geometry and between \(0.8894\) and \(0.9468\) across the three pre-specified non-control geometries. The corresponding oracle-static point gaps are \(0.0666\), \(0.0916\), \(0.0388\), and \(0.0342\).

The invariant oracle result across V1b geometries is a useful structural check: the clean reference and future stream are held fixed, and a correct oracle removes the contaminated sorted positions before applying the same clean-rank construction. Geometry affects the robust and confidence-band procedures, but not the underlying canonical clean rank sequence.

We do not derive an analytical efficiency bound for the oracle-static gap. Characterizing its dependence on \(K\), \(m\), horizon, and alternative strength remains open.

\subsection{What V1b changes}

The geometry stress test rules out a simple explanation of the original power gap: that both contaminants at value \(2\) created an unusually favorable setting for static coupling. The fixed confidence-band comparator improves substantially when contaminants are moved into the support or split across tails, rising from \(0.1152\) to as high as \(0.3948\) at the primary point. The static method nevertheless remains between \(0.8894\) and \(0.9468\).

This does not establish uniform dominance over all geometries. It does show that the main empirical finding survives a pre-specified set of materially different static/exogenous contamination arrangements without changing the detector.

\subsection{Why the band comparison is conservative}

The confidence-band baselines represent the contaminated reference through a simultaneous CDF uncertainty set. This remains valid without identifying contaminated positions, but the resulting uncertainty set can be substantially larger than the set compatible with one persistent contamination pattern. The final development audit intentionally favors the band family by sweeping a frozen parameter grid and reporting the best configuration ex post for each condition.

Because the ex-post envelope is not one precommitted test, it is used only as a descriptive adversarial benchmark. V1b fixes \texttt{band\_mix}, \(\delta=0.045\), before its geometry outcomes and reports paired Monte Carlo uncertainty against that comparator. Because the comparator was informed by the earlier Publication V1 analysis and the V1b control reuses those paths, these intervals are not presented as selection-adjusted confirmatory inference.

\subsection{Marginal versus conditional guarantees}

Theorem~1 inherits the marginal fixed-reference interpretation of predictive-rank martingales \cite{KuangXia2026}. It averages over the random clean reference sample. CCTM instead targets a calibration-conditional guarantee with high probability over a clean reference draw \cite{ShaerEtAl2026}. These are different inferential targets. We do not claim that one uniformly dominates the other.

\subsection{Connection to CFAR and receiver data}

The radar motivation is direct but limited. Classical CFAR uses neighboring reference cells to set a threshold for a cell under test, with order-statistic and censored variants addressing nonhomogeneous windows and multiple interferers \cite{Rohling1983}. Our setting instead reuses one finite reference bank repeatedly, permits persistent contaminated cells, and requires optional-stopping-safe sequential evidence.

The semi-synthetic RTL-SDR experiments provide an engineering bridge rather than a theorem check. Real receiver/background IQ preserves quantization and environmental structure, while controlled post-ADC injection preserves known ground truth. V2 exposes a non-saturated transition: at \(-18\) dB static crosses on 14/30 captures versus 0/30 for the fixed band comparator, and at \(-15\) dB the corresponding counts are 22/30 versus 12/30. These observations are qualitatively consistent with the simulation advantage, but serial dependence prevents interpretation as hardware type-I validation.

\subsection{Structural versus distributional robustness}

The model illustrates a broader distinction. Distributional robustness enlarges the set of possible laws. Structural robustness can exploit restrictions on how nuisance variables persist. Here the nuisance parameter is not a new arbitrary corruption at every time; it is the fixed set \(C^\star\) shared by all future comparisons. The theorem itself is simple once that structure is stated, but recognizing and preserving this constraint is what makes the construction useful.

\subsection{Computational considerations}

Exact enumeration requires
\[
  \binom{K}{m}
\]
candidate processes. This is practical for the frozen experiments with \(m=2\), but does not scale well when both \(K\) and \(m\) increase. The contiguous-partition dynamic program solves the fixed-categorical case in \(O(Km^2)\) time and \(O(Km)\) memory. The full frozen convex portfolio lacks the same additive structure, so no polynomial-time exact algorithm is claimed for it.

\subsection{Finite-reference information ceiling}

For fixed \(n\),
\[
  P\sim\mathrm{Dirichlet}(1,\ldots,1),
  \qquad
  R_t\mid P\overset{\mathrm{iid}}{\sim}\mathrm{Categorical}(P).
\]
As the online horizon grows, the stream increasingly reveals the realized spacing vector \(P\), but does not eliminate randomness from the finite reference itself. Some persistent rank patterns under alternatives can therefore remain compatible with rare null spacing realizations. This finite-reference ceiling is consistent with the limitation identified for clean PRMs by Kuang and Xia \cite{KuangXia2026} and motivates our focus on finite-reference power rather than universal fixed-\(n\) consistency.
